\section{K-023: cyclic tensor trace as a hook-cubic invariant series}

Let $W=\mathbb C^n$ and put
\[
 \mathcal U_n=W\oplus\mathbb S_{(2,1)}W.
\]
Give the two summands central weights one and three, respectively.
For the coordinate algebra of $\mathcal U_n$, use this positive
weighted polynomial grading.  Every $SL_n$ invariant has weighted
degree divisible by $n$; let $b_n(s)$ be its dimension in weighted
degree $ns$, and put $B_n(t)=\sum_{s\ge0}b_n(s)t^s$.
Let $\tau$ cyclically permute the three tensor factors in the
triple-tensor quotient ring $R_n$.

\begin{theorem}[K-023]
For every $n\ge2$ and $s\ge0$,
\[
 \operatorname{tr}(\tau\mid (R_n)_s)=b_n(s),\qquad
                         k_n(s)\equiv b_n(s)\pmod3.
                                                               \tag{K-023a}
\]
In fact $k_n(s)-b_n(s)$ is a nonnegative integer multiple of three.
Thus the two integral Hilbert series satisfy
\[
                         H_n^{\rm ten}(t)\equiv B_n(t)\pmod3
                           \quad\hbox{in }\mathbb F_3[[t]].
                                                               \tag{K-023b}
\]
\end{theorem}

\begin{proof}
We can perform the character computation with polynomial
representations on $W$ instead of $W^*$: the simultaneous
contragredient automorphism of $SL_n$ preserves the dimensions
of invariant spaces and commutes with permutation of the tensor
factors.  In weighted degree $d=ns$, Schur--Weyl duality gives
\[
 (R_n)_s\simeq
       ([\lambda]\otimes[\lambda]\otimes[\lambda])^{S_d},
       \qquad \lambda=(s^n),                              \tag{K-023c}
\]
compatibly with the cyclic factor permutation.  Indeed the
symmetric power of $W\otimes W\otimes W$ is the diagonal
$S_d$-invariant subspace in $(W^{\otimes d})^{\otimes3}$, and
$(W^{\otimes d})^{SL_n}$ is the Specht module $[\lambda]$.
For $d=0$ all these spaces are interpreted as the trivial line.

If $A$ is an endomorphism of a finite-dimensional vector space,
then the trace of cyclic permutation composed with $A^{\otimes3}$
is $\operatorname{tr}(A^3)$; expansion in a basis proves this by
contracting the three cyclic indices.  Averaging the diagonal
$S_d$ action in (K-023c) therefore gives
\[
 \operatorname{tr}(\tau\mid(R_n)_s)
   =\frac1{d!}\sum_{\pi\in S_d}\chi_\lambda(\pi^3).
                                                               \tag{K-023d}
\]

We compute these character averages simultaneously.  Write $p_j$
for the power-sum symmetric functions and $z$ for weighted degree.
The Frobenius character identity
$p_\mu=\sum_\lambda\chi_\lambda(\mu)s_\lambda$ shows that the
coefficient of $s_\lambda z^d$ in
\[
         \sum_{d\ge0}\frac{z^d}{d!}
                        \sum_{\pi\in S_d}p_{\operatorname{type}(\pi^3)}
                                                               \tag{K-023e}
\]
is precisely the average in (K-023d).  A cycle of length $j$
remains a $j$-cycle under cubing when $3\nmid j$, whereas a
cycle of length $3r$ becomes three $r$-cycles.  The exponential
formula for permutations consequently turns (K-023e) into
\[
 \exp\left(\sum_{3\nmid j}\frac{z^jp_j}{j}
                   +\sum_{r\ge1}\frac{z^{3r}p_r^3}{3r}\right).
                                                               \tag{K-023f}
\]
The identity
\[
                      s_{(2,1)}=\frac{p_1^3-p_3}{3}
\]
now gives the exact equality
\[
\begin{split}
 &(K\text{-}023f)\\
 &\quad=\exp\left(\sum_{r\ge1}\frac{z^rp_r}{r}
              +\sum_{r\ge1}\frac{z^{3r}}{r}
                         p_r[s_{(2,1)}]\right).
\end{split}                                                   \tag{K-023g}
\]
To check it directly, use
$p_r[s_{(2,1)}]=(p_r^3-p_{3r})/3$; the negative terms cancel
exactly the $j$ divisible by three in the first sum.
Expression (K-023g) is the character of the symmetric algebra
on $W\oplus\mathbb S_{(2,1)}W$, with weights one and three.
Specialization to $n$ variables and extraction of the coefficient
of $s_{(s^n)}$ therefore gives $b_n(s)$.  This proves the trace
identity.

Finally $(R_n)_s$ and its factor permutation are defined over
$\mathbb Q$.  An operator of order three on a rational vector
space has characteristic polynomial
\[
                         (X-1)^a(X^2+X+1)^b
\]
for nonnegative integers $a,b$.  Its dimension is $a+2b$ and
its trace is $a-b$.  Their difference is $3b$.  Applying this
to $\tau$ proves the remaining assertions.
\end{proof}

\paragraph{Representation-ring interpretation.}
The proof is an identity in the completed $\lambda$-ring of
polynomial representations.  Cubing a permutation implements
the third Adams operation on its character; the cycle-index
exponential becomes the symmetric algebra on the actual Schur
module $W\oplus\mathbb S_{(2,1)}W$.  The reduction modulo three
is the ordinary integral character identity for a rational
$C_3$ representation.  No inference from a virtual positive
class to geometric purity is used.

The symmetric-function identities in the proof are the standard
Frobenius and Schur-operation conventions of Macdonald,
\emph{Symmetric Functions and Hall Polynomials}, second edition,
Chapter I, \S7; the calculation above supplies the particular
cycle-index specialization needed here.  No historical priority
is claimed for the general higher-indicator identity.

The general root-number character identity is due to Scharf (1991),
\emph{Die Wurzelanzahlfunktion in symmetrischen Gruppen},
\emph{Journal of Algebra}139,446--457.  A primary modern statement
is Adin--Heged\H us--Roichman, arXiv:2312.08904v2,
Definition2.1 and Theorem2.2, printed p.5; the latter cites Scharf.
The proof above specializes that known identity to cubic cycles and
then to the rectangular invariant spaces, making the trace and
regrading conventions explicit.
